\subsection{Local descent on an adapted pair of theta towers}
\label{tcr:main}

Write \(d=a+a'\), \(k=(a-a')/2\), and \(s=p+q\). In the
\(A\)-step take \(q=0<p\), remove a plane from \(V,W\), and put
\(z=k,\ \xi=\alpha|\cdot|_E^{a'/2}\). In the \(B\)-step take
\(q>0\), remove a plane from \(V',W\), and put
\(z=k+1,\ \xi=\beta|\cdot|_E^{a/2}\).
The one-plane evaluation \(e_0\) belongs to the unchanged source
and the two-plane cone operator \(T\) to the changed source. Subscripts
on these maps denote their local factors. The adaptedness assumptions
are those of \S\ref{newtheta:global}: the two integral lattices
split into their two lines and their orthogonal complements, the
finite designated bases have the prescribed columns, and the Cartan
involutions pair the two lines and restrict to the lower involutions.

\begin{lemma}[Square integrability and kernels]\label{tcr:kernels}
Each upper polynomial period, with the stipulated primed support,
belongs to \(L^2([U(W)])\). For any of the matching maps, the period
on each side factors through the restricted tensor product of its
small local theta quotients. A nonzero such period map is injective
on that tensor product.
\end{lemma}
\begin{proof}
The periods are smooth, of uniform moderate growth, finite under the
enveloping center and compact-finite at infinity, as in the proof
of Proposition~\ref{threal:detectline}. Apply the cuspidal-exponent
criterion for square integrability, with the constant term divided
by the square root of the modulus of its parabolic
\cite[I.4.11]{moeglinWaldspurger}. The center of the unitary group
has compact adelic quotient.

Consider a cuspidal exponent on a Levi of a flag of isotropic
subspaces. If the first subspace has dimension \(r>1\), its first
Levi block is \(\mathrm{GL}_r(\mathbb A_E)\). Any cuspidal projection
there is generic and unitary after a central real twist. Here a
cuspidal projection is computed by pairing the flag constant term
on \(M(F)\backslash M(\mathbb A_F)^1\), where all rational
character moduli are one, with a rapidly decreasing cuspidal test
on the Levi \(M\). Uniform moderate growth suffices for convergence;
the central height variables remain free for extracting exponents.
Fix the other cuspidal factors and vary the first cusp vector. A
nonzero pairing gives an equivariant functional with values in
the smooth contragredient of that cuspidal factor, using finite
compact invariance of the local test. An action off the norm-one
subgroup is accounted for by a compensating split-central
translation, hence by the extracted central exponent.
At a completely
split unramified prime with \(\nu_w(\pi)\) of infinite order it has
to be a subquotient of the first block in the Jacquet module of the
upper small quotient. Indeed the full flag constant term can be taken
first along the \(r\)-plane and then in its remaining unitary factor.
Integration there and cuspidal projection commute with the first
block action. The local calculation in Proposition~\ref{threal:detectline}
rules out precisely these subquotients for \(r>1\). This argument
uses the full split source-coinvariant calculation at that one
prime and does not require the period to be in \(L^2\).

If \(r=1\), the two line components are smaller periods by
Propositions~\ref{threal:coneA}, \ref{tpd:unfold} and
Lemma~\ref{tpd:so}, including their zero-orbit formulas. A missing
source plane gives no corresponding nonzero orbit. The same
calculation applies to the remaining flag on each contributing
smaller period. For this calculation rational hyperbolic coordinates
suffice: at the places where the resulting bases are not integral
they change finitely many local tests. At infinity the estimates of
Proposition~\ref{threal:decay} apply to Schwartz tests and their
derivatives. Center-finiteness of the successive coefficient as an
automorphic function follows also by Harish-Chandra projection from
the preceding coefficient. Thus every general linear block of a
contributing cuspidal exponent is a line. If there are \(j\) such
blocks, write \(T_i=|t_i|_E\) for their parameters in flag order.
The real exponents \(b_i\) of the unnormalized term satisfy
\(b_i\le a/2\), since at every stage they are half the rank of one
of the two current sources. The square root of the modulus has
exponents \((d-2i+1)/2\): the line radical for a space of dimension
\(d-2i+2\) has determinant modulus \(T_i^{d-2i+1}\) (its middle
module has dimension \(d-2i\) over \(E\), and its central
\(F\)-line has modulus \(T_i\)). Hence
\[
        b_i-(d-2i+1)/2\ \le\ -(a'-2i+1)/2<0
        \qquad(1\le i\le j\le s).
\]
Every nonzero generating ray of the closed positive chamber
\(\log T_1\ge\cdots\ge\log T_j\ge0\) has strictly negative
normalized exponent. The criterion gives the assertion.

Fix the other local tests. Source invariance factors the resulting
map into \(L^2\) through the big theta module of the trivial
representation, which has finite length and a unique simple quotient
with multiplicity one in its cosocle
\cite{howeTranscending,ganTakedaHowe,minguezHowe}. An admissible
finite-length image in a unitary representation is semisimple.
For a submodule take its Hilbert closure: the projection to a
compact-open fixed space at a finite place or a compact isotypic
space at a real place shows that the closure intersects the
algebraic image in the submodule. These spaces are finite-dimensional,
so orthogonal complements taken on them give an algebraic splitting;
at infinity it is stable under derivatives by skew adjointness on
smooth vectors. Thus each local map kills the kernel of the prescribed
small quotient. At the support place \(v_1\) on the primed side one
uses the supported coinvariants, which map isomorphically to this
same small quotient by Proposition~\ref{split:kernel}. Applying
the universal property of a tensor product to finite sums and then
enlarging the finite set of places gives the factorization.

The kernel of the map on the resulting restricted product is stable
under independent finite-place translations and the real Lie and
compact actions. A submodule containing a nonzero vector contains
the full finite tensor product supporting that vector: the individual
simple factors have scalar endomorphisms, and the algebraic density
theorem gives simplicity of their finite external product. Enlarging
the finite set, using the nonzero spherical vectors at the new
places, gives the whole restricted product. This proves injectivity
whenever the period map is nonzero.
\end{proof}

\begin{lemma}[Nonzero smaller images at infinity]\label{tcr:bump}
Suppose the smaller period on the changed-source side is nonzero
on polynomial tests. Then \(P_{1,\infty}T_\infty\) is nonzero on
upper polynomial tests, where \(P_{1,\infty}\) is its real matching
map. For the other side \(P'_{1,\infty}e_{0,\infty}\) is nonzero
whenever its smaller real quotient is nonzero.
\end{lemma}
\begin{proof}
Choose smaller finite tests and a target argument at which the
smaller period is a nonzero functional \(D\) of the smaller real
test. It is continuous for the Schwartz topology by
Corollary~\ref{threal:limits} and source-invariant. Its restriction
to polynomial tests factors through the smaller real quotient by
Lemma~\ref{tcr:kernels}. For \(e_0\), multiply any smaller Fock
polynomial by the vacuum on the complementary summand, whose
evaluation is nonzero. The positive complex structure preserves
this decomposition by adaptedness. Thus the one-plane evaluation
is onto on polynomials.

For \(T\), use the mixed coordinates with modulation space
\(Y\otimes Ex'\oplus Ex\otimes W_1\), where \(Y\) is the changed
source. Before the second evaluation the first position is
\(v\in Y_\infty\). For a real compact source element \(h\) denote
the evaluation at \(v=u x\), followed by the second evaluation
and \(D\), by \(A_h(u)\), where \(u\in\prod_{w\mid\infty}\mathbb C\).
In these coordinates multiplication of the initial test by a
scalar test \(b(v)\) multiplies \(A_h(u)\) by \(b(h^{-1}ux)\).
Indeed the source acts on the first position by pullback and on
the residual oscillator by its Weil action. With \(h=1,u=1\),
the evaluation onto the smaller Schwartz space is surjective
by extension in the omitted linear coordinates.

In the \(A\)-step the scalar under the integral for \(DT_\infty\)
is \(A_h(u)\) times the everywhere nonzero Mellin density on
\((\mathbb C^\times)^{[F:\mathbb Q]}\). In the \(B\)-step one
instead pairs the one-coordinate Fourier transform of \(A_h\)
at each place with the Mellin density of exponent \(k+1\), by
the partial Fourier identity in Proposition~\ref{tpd:unfold}.
We describe the transformed density on \(u\ne0\). Relative to
additive area on \(\mathbb C\), the original density is a nonzero
constant times
\[
          (l/|l|)^{m_w}|l|^{2k},\qquad |m_w|=m_{\alpha,w}-m_{\beta,w}>d .
\]
Complex conjugation of the displacement or reversal of the Fourier
pairing can reverse the angular weight. The Fourier transform of
this tempered distribution has a smooth nowhere-zero density on
\(\mathbb C^\times\). Indeed it is equivariant under rotations and
positive dilations. On their free transitive orbit such a distribution
is a scalar times the corresponding character times Haar density:
divide by the character and use translation invariance in angle and
logarithmic radius. If its scalar vanished it would be supported
at the origin, hence be a finite sum of derivatives of the delta
distribution. Its inverse Fourier transform would then be a
polynomial of homogeneous degree \(2k\) and angular weight of
absolute value \(>2k\), which is impossible by expanding in
\(l,\bar l\). This argument applies separately at every real place.

Here is why integration over the compact directions cannot cancel
this density. In the first mixed coordinates
\((h,u)\mapsto h^{-1}ux\) maps the compact group times the
nonzero scalars onto the product of punctured null cones. At each
place the two compact unitary groups are transitive on the two
unit spheres, and the lengths of the two components of a null
vector are equal. The fibers are the compact line stabilizers,
acting by \((h,u)\mapsto(jh,\eta u)\) for \(jx=\eta x\).
The factor acting on the smaller oscillator is the smaller compact
source and disappears after \(D\). The scalar on this stabilizer
also cancels: for the unchanged radial measure the source determinant
factor on the line is \(\nu^{-1}(\eta)\) in the \(A\)-step and
\(\nu(\eta)\) in the \(B\)-step, as in the source Levi calculations
of Proposition~\ref{threal:coneA} and Lemma~\ref{tpd:mellin}.
In the first chart the weight on \(u\) is \(\nu(u)\) in the
\(A\)-step. In the other step the Fourier transform carries a
rotation of \(u\) to its dual rotation, so the transformed angular
weight is \(\nu^{-1}(u)\). Explicitly the \(B\)-chart functions
after applying \(D\) satisfy
\[
        B_{jh}(l)=\nu(\eta)B_h(\eta l),\qquad |\eta|=1.
\]
The argument \(\eta l\) follows because the displacement is
\(\bar l\,y\otimes x'\): the Fourier kernel pairs \(u\) with \(l\)
by a real trace of a fixed nonzero multiple of \(ul\).
Substitution in the defining Mellin integral with weight \(\nu(l)\)
cancels the factor \(\nu(\eta)\). This proves invariance of the integrated density
along the compact fibers in either chart.

Choose first a test with \(A_1(1)\ne0\), then multiply it by a
nonnegative bump \(b(v)\) supported sufficiently near \(x\) on
each punctured cone and positive at \(x\). The pullback of that
support has all \(|u_w|\) bounded away from zero and infinity.
The Fourier distribution may therefore be evaluated there using
its smooth density. After compact-fiber integration the remaining
integrand, apart from \(b\), is continuous and nonzero at \(x\);
in a small neighborhood a fixed phase makes its real part positive.
This proves \(DT_\infty\ne0\) on Schwartz tests. Compact averaging,
Fourier transformation, restriction and the absolutely convergent
Mellin moment are continuous in the Schwartz seminorms, as in
Lemma~\ref{tpd:mellin}. Hermite expansions approximate any
Schwartz test by polynomials times the vacuum in these seminorms.
The output for an upper polynomial is again lower polynomial:
decompose it over the smaller positive summand and its orthogonal
complement; the compact orbit spans a finite-dimensional polynomial
space and all evaluation and displacement act in the complement.
Consequently some upper polynomial has \(DT_\infty\ne0\), as claimed.
\end{proof}

\begin{proposition}[Descending the test maps]\label{tcr:descend}
Suppose both smaller period maps are nonzero. After application of
the smaller matching map, each of \(T_v,e_{0,v}\) factors through
the upper matching map, at every inert place and at the product of
the real places (on polynomial vectors there). Both resulting maps
intertwine the target line parabolic with
\(\xi\boxtimes\pi_1\) inflated trivially on its radical.
At infinity this means its Lie algebra and the compact part of its
Levi. The two real maps so obtained are nonzero.
\end{proposition}
\begin{proof}
First produce auxiliary upper tests at every other place for each
of the two operations, such that their smaller images are nonzero.
At infinity use Lemma~\ref{tcr:bump}. At a ramified place use the
standard matched pair of Proposition~\ref{tpd:ram} or
Corollary~\ref{threal:ramA}. At an unramified split place use the
vacuum in Proposition~\ref{tpd:ordinary}. At \(v_1\), if a
supported primed test is required, use the surjection of supported
tests onto the whole upper quotient. All the indicated local maps
at that split place already factor through the upper quotient
there by those propositions (one may compare with the zero test on
the opposite side). In the \(A\)-step the smaller unchanged primed
source is compact at infinity and needs no support for integration.
When the smaller supported space is required for matching, one
can also use the representatives of Lemma~\ref{tpd:so}: the same
matrix proof of its lifting assertion works for \(q=0\).
In the \(B\)-step the changed-source operator itself preserves the
full-rank condition by Lemma~\ref{tpd:mellin}.

At an inert place the following vacuum calculation requires no
factorization. The tensor lattice at the place is self-dual for
the additive-character symplectic pairing. The prime is nondyadic.
In the two-plane symplectic decomposition its summands in the
removed tensor blocks have the following description. Each isotropic
line tensor a residual space pairs perfectly with the opposite
isotropic line tensor that space. The four products of the two pairs
of lines split into two such dual pairs. Adaptedness and self-duality
make these pairs of lattices mutually dual, including any fractional
line ideals. Take the removed modulation space and these opposite
summands as its dual position space. In those positions the basic
vector is the lattice indicator times the smaller basic vector,
up to scalar: invariance under the modulation lattice imposes
support on its position dual and invariance under the full dual
translation lattice forces constancy there. The tensor factor on
the residual space is itself the smaller basic line by the same
Heisenberg characterization. This argument also applies to the
one-plane evaluation. Normalize each evaluation on the vacuum.
The source compact fixes the basic vector, by the unramified
oscillator vacuum calculation in Proposition~\ref{threal:allinert}.
The displacement is a single position coordinate in the respective
chart. Thus shifted evaluation gives precisely the smaller basic
vector times \(1_{\mathfrak a_v}(l)\), where the integral translation
ideal \(\mathfrak a_v\) is the localization of
\[
 \mathfrak b\,\mathfrak D_E^{-1}\bar{\mathfrak c}^{-1}
       \quad(A),\qquad \mathfrak b^{-1}\bar{\mathfrak c}\quad(B).
\]
With multiplicative unit volume one its Mellin integral is
\[
  u_v=\nu_v(b_v)|b_v|_{E_v}^{z}
           (1-\nu_v(\pi_v)|\pi_v|_{E_v}^{z})^{-1},
        \qquad b_v\mathcal O_{E_v}=\mathfrak a_v .
\]
This is nonzero, including when \(\nu_v(\pi_v)=1\). All but
finitely many \(b_v\) are units, and the product of these basic
factors converges absolutely to a nonzero number since \(z>1\).

Now put a vector in the kernel of an upper matching map at the
place under consideration and the chosen auxiliary tests everywhere
else. Lemma~\ref{tcr:kernels} kills its upper period identically.
In particular its \(\xi\) coefficient is zero. This coefficient
is a fixed nonzero Haar multiple of the smaller period on the
tensor of the transformed tests, by the constant-term formulas.
For the primed cone formula for a general upper test, first take
the Levi difference of Proposition~\ref{tpd:unfold}; on the
\(\beta\)-coefficient it acts by the difference of the two
characters at the chosen idele, a nonzero scalar, and on \(T\)
it acts by that same scalar. For a zero evaluation with supported
representative, make the correction at \(v_1\) before reading the
other local factor. By smaller injectivity
(Lemma~\ref{tcr:kernels}) the tensor of lower images vanishes.
All its other factors are nonzero, so the one at the place under
consideration vanishes. This proves the claimed factorizations.

The target equivariance can be read from the same formulas. For
a local element of the line Levi or radical take its right translate
in the upper constant term. At the identity the radical acts
trivially by left invariance, the line acts by \(\xi\), and the
smaller target group translates the smaller period. Subtract the
corresponding two smaller expressions and again cancel the fixed
nonzero auxiliary factors by smaller injectivity. At infinity
one takes compact Levi translates and differentiates the identity
for one-parameter subgroups at zero; these operations preserve the
polynomial domain in the required differentiable sense. Continuity
for passing the derivative through the periods follows from
Proposition~\ref{threal:decay}. For a one-parameter element outside
the compact the identity before differentiation holds for Schwartz
tests by the same absolutely convergent integrals. The nonzero
real images follow from Lemma~\ref{tcr:bump}.
\end{proof}

\begin{proposition}[Inert and real uniqueness]\label{tcr:unique}
Under the assumptions of Proposition~\ref{tcr:descend}, the two
descended maps at an inert place satisfy \(D_{T,v}=u_vD_{0,v}\).
At infinity \(D_{T,\infty}=c_\infty D_{0,\infty}\) for a unique
nonzero scalar. For fixed real models, upper and lower real
matching maps and measures, this number is independent of every
finite character, matching quotient and choice of admissible
Weil actions in the two towers.
\end{proposition}
\begin{proof}
At an inert place the upper spherical vector generates its
irreducible representation. The line parabolic and the lattice
stabilizer satisfy \(G=PK\) by the integral Witt argument in
Proposition~\ref{threal:coneA}. Hence the spherical vector generates
even under \(P\). The comparison of its two images in the preceding
proof gives the identity on the whole representation.

Here is the compact weight calculation at one real place.
Write \(R=d-s,S=s,P=a-p,Q=p,A=m_\alpha,B=m_\beta\), and
\(\Delta=A-B>d\). Put
\[
 c_\pm=\tfrac12(B\pm(P-Q)),\qquad
 h_+=c_++\tfrac12(\Delta-R+S),\quad
 h_-=c_-+\tfrac12(\Delta+R-S).
\]
The bottom \(U(R)\times U(S)\)-type of
Proposition~\ref{threal:allreal} has weights
\[
 ((h_+)^P,(c_+)^{R-P};(h_-)^Q,(c_-)^{S-Q}),
\]
where superscripts denote repetitions. To see the normalization,
start from the two target compact vacuum weights \(\pm(P-Q)/2\);
the source determinant in the joint harmonics adds
\((\Delta-R+S)/2\) in the first \(P\) positions and
\((\Delta+R-S)/2\) in the first \(Q\) positions, and the target
splitting adds \(B/2\) throughout.
Alternatively the bottom-layer formula for the ordered two-block
parabolic is its inducing differential plus
\(\rho(\mathfrak u\cap\mathfrak p)-\rho(\mathfrak u\cap\mathfrak k)\)
in Paul--Trapa's cover normalization: in the first positive block
this correction is \(((S-Q)-(R-P))/2\), in the first negative
block its negative, and in the remaining positive and negative
blocks it is respectively \((P-Q)/2\) and its negative.
The inducing parameter is in the good range by \(\Delta>d\);
the bottom type occurs once by the bottom-layer formula for
cohomological induction (the inducing Levi type is a character).
In ordinary, uncovered normalization one subtracts \(\rho(\mathfrak u)\)
from the inducing differential. The two character exponents then
become \((A-a')/2,(B+a)/2\); their difference is
\((\Delta-d)/2>0\), so the bottom-type theorem applies.
\cite[\S2 and Theorem~4.1]{paulTrapaHowe}
\cite[Theorem~5.3 and Proposition~6.1]{voganZuckerman}

The compact Levi on an isotropic line acts by the same circle
\(\eta\) in one positive and one negative coordinate, and by
\(U(R-1)\times U(S-1)\) on the other coordinates. Indeed a null
vector is a sum of two vectors of equal norm, and an element
carrying it to \(\eta\) times itself acts by \(\eta\) on each.
For a \(U(n)\) weight \(\lambda_1\ge\cdots\ge\lambda_n\) the
\(U(n-1)\)-weights interlace it with multiplicity one, and their
omitted circle weight is \(\sum\lambda_i-\sum\mu_i\). Its maximum
\(\lambda_1\) forces \(\mu_i=\lambda_{i+1}\) for all \(i\);
its minimum \(\lambda_n\) forces \(\mu_i=\lambda_i\).
In the \(A\)-step both high blocks occur and \(h_++h_-=A\) is
the maximum diagonal-circle weight. In the \(B\)-step both low
blocks occur and \(c_++c_-=B\) is the minimum. These are exactly
the respective angular characters of \(\xi\). The extreme block
in the upper bottom type thus consists of a single irreducible
type of the smaller compact, with one entry deleted in each sign.
It is the lower bottom type by the same weights formula:
\(P-Q\) and \(R-S\) do not change. It occurs once in \(\pi_1\).
The compact restrictions of the two parabolic maps on the upper
bottom type are therefore proportional. The Iwasawa Lie algebra
identity \(\mathfrak g=\mathfrak p+\mathfrak k\) and the
Poincar\'e--Birkhoff--Witt theorem show that \(U(\mathfrak p_{\mathbb C})\)
acting on this compact type generates the entire irreducible module.
For the product of real places take the exterior tensor product
of the bottom types and apply the same argument componentwise.
This proves dimension at most one for the required space of maps.
Both are nonzero, so the scalar is unique and nonzero.

For fixed base data the infinite characters are fixed by their
specified exponents \(m_\alpha,m_\beta\). The doubled Siegel law
fixes the resulting real oscillator actions: the ratio of two
Heisenberg implementations with the same doubled Levi law is a
character trivial on that Levi and on the root unipotents, hence
trivial on the doubled group, as in Proposition~\ref{threal:inert}.
The compact integral defining \(T_\infty\) and its Mellin weight
therefore agree for all the allowed systems. The fixed realizations
of the upper and lower real maps give the same proportionality
equation on every polynomially matched real pair in all these
systems. Since some real test has nonzero zero-evaluation image
and its upper image can be matched on the other side, this equation
determines the same number \(c_\infty\) for all systems. If
there are no such systems the uniform comparison is vacuous.
\end{proof}

\begin{theorem}[The scalar for the two undoubled periods]\label{tcr:global}
For every triple of \S\ref{newtheta:global} and every allowed set
of real and finite matching maps, the primed period is nonzero
on some supported test, and there is a unique \(c\ne0\) such that
\(I(\Phi)=cI'(\Phi')\) on all matched pairs. For given adapted upper
and lower geometries and given upper and lower real matching maps,
there is \(C\ne0\) independent of all admissible \(X\), Weil actions
and finite matching maps such that
\[
 \begin{array}{ll}
 c=C\,F_X(1-k)L^{S_0}(\nu,k)c_1,
       &(q=0<p;\ e_1=e_2=+1,\ s_{\rm unit}=+1),\\[2pt]
 c\,L^{S_0}(\nu,k+1)=c_1/(C F_X(-k)),
       &(q>0;\ e_1=e_2=-1,\ s_{\rm unit}=+1).
 \end{array}
\]
Here \(F_X\) is the finite factor of \S\ref{newtheta:global},
and \(c_1\) is the scalar for the specified smaller matching maps.
\end{theorem}
\begin{proof}
First work on an adapted pair of geometries and suppose the
assertions of nonvanishing and proportionality hold in the smaller
rank. In particular both smaller period maps are nonzero: supported
primed vectors are matched by surjectivity of the opposite quotient
and the distinguished lower full-rank quotient. Use polynomial
inputs initially; nonzero polynomial inputs are available by
Corollary~\ref{threal:limits} and density. For \(T\) use the
integral with multiplicative units of volume one and local source
compacts of volume one. Normalize finite mixed evaluations on
their vacua. Rational coherent evaluations differ from this system
by fixed scalars with almost every scalar one
(Proposition~\ref{threal:coneA}). At infinity keep the real
evaluations fixed.

On upper matched tests, the restricted tensor product of lower
images of \(T\) is a scalar times that of \(e_0\). At the real
places use Proposition~\ref{tcr:unique}; at the finite places use
Propositions~\ref{tpd:ordinary}, \ref{tpd:ram},
Corollary~\ref{threal:ramA} and Proposition~\ref{tcr:unique}.
For a finite tensor sum this follows by applying the tensor product
of the descended linear maps to the equality of upper tensors in
the definition of polynomial matching. Outside the finite test set
one must retain the product of the vacuum multipliers: these
converge absolutely and have nonzero product. The translation
ideal contributes \(N(\mathfrak a)^{-z}\nu(\mathfrak a)\).
It is a unit at the ramified factors. As computed after
Lemma~\ref{tpd:mellin} and Proposition~\ref{tpd:ram}, its
character value is \(\nu(\mathfrak b)\nu(\mathfrak c)\) in \(A\)
and its inverse in \(B\). The partial Euler factors at the primes
of \(Q\) are one since \(\nu\) is ramified there. The ramified
factors are the stated Gauss factors, with arguments \(1-k,-k\)
respectively and with unit sign \(+1\) in both cases.
All these factors are nonzero. The remaining factors consist of
\(c_\infty\), the absolute ideal norm power, the fixed changes of
evaluation, and the quotient of the two Haar normalizations in
the orbit formulas. They give a single \(C\ne0\) for the stated
fixed data by Proposition~\ref{tcr:unique}. The Haar measures
are part of the models; no equality of Haar masses in different
dimensions is needed.

Use the smaller identity on these two tensors. In \(A\), the
unprimed \(\xi\)-coefficient is the cone and the primed one is the
zero evaluation, giving the first formula. In \(B\), the primed
coefficient is the cone and the unprimed one is the zero evaluation,
giving the reciprocal in the second formula.
Propositions~\ref{threal:endA} and~\ref{threal:endB} show that
the resulting scalar compares the upper periods at every target
argument. The zero-coefficient side is nonzero for some upper
test: choose its real polynomial and finite factors as in
Proposition~\ref{tcr:descend}, then apply smaller injectivity to
their nonzero tensor. Choose a matching vector on the other side,
using the supported surjection at \(v_1\) when the other side is
primed. The comparison just proved gives nonzero upper periods
on both tests. This also proves uniqueness of \(c\). Passage to
every matched limit is Corollary~\ref{threal:limits}.

To apply this argument by induction on \(p+q\), start at
Theorem~\ref{threal:base}. For the induction step the integral
plane construction of Proposition~\ref{famtheta:plane} supplies
the hyperbolic summands with their ideal lattices in the changing
source and target. On a given rational hyperbolic plane choose
Cartan involutions interchanging the two null lines at infinity
and an arbitrary compatible Cartan on its complement. At the
designated split places take the new adapted integral bases in the
prescribed column order. These choices do not change the rational
spaces, the lattices, the Haar measures or their theta integrals.
The smaller tensor has a rational symplectic basis by symplectic
elimination, giving its usual rational Schr\"odinger model.
Its basic vector is the lattice Heisenberg eigenline. It can also
be obtained by partial evaluation of the upper basic vector in the
adapted coordinates above: the removed mutually dual lattices
give the position indicator in those coordinates, leaving the
smaller eigenline. At a split dyadic place one uses the two isotropic
idempotent halves as in Proposition~\ref{tpd:ordinary}; their
individual quadratic phases vanish. In rational coordinates the
basic lattice gives the ordinary integral indicator almost everywhere.
Use any fixed lower source Haar measures. For the two lower
oscillators use the same pairs of splitting characters
\((\beta,\beta/\alpha)\), \((\alpha,\alpha/\beta)\).
Subtracting two from either rank preserves the quadratic character
restriction required for each unitary splitting. Use the local
splitting and its doubling law
\cite[(1.16) and Proposition~2.2]{harrisKudlaSweet}, and their
adelic product for the two Hecke characters
\cite[\S\S2--3]{wuThetaTower}. On rational elements the character
factors, absolute determinant factors and Weil indices have product
one, giving the rational theta action through Poisson summation.
Equivalently the doubling in a rational diagonal polarization
uses the prescribed determinant character on its Siegel Levi
and evaluation at zero; restricting to the two isometric summands
recovers the chosen single Weil action. All smaller gap inequalities
hold since \(p+q\) decreases by one, \(a'\) decreases by at most
two and \(k\) decreases by at most one. The local matching results
in \S\ref{threal:main} apply at both levels. The induction
hypothesis and the preceding adapted argument now apply.

We finally transport the existence and uniqueness assertion to the
original markings. At infinity choose separately in each source
and commonly in the target real isometries conjugating the original
Cartan involutions to the adapted ones. Their oscillator operators
carry the Gaussian line and its creation derivatives onto the new
ones by Heisenberg covariance. On the period the source translation
disappears by Haar invariance and the common target translation
right translates both functions by the same element. Under this
transport the original and new pairs of real simple quotient maps
differ, after identifying their targets, by one nonzero scalar
for each side, by real Howe uniqueness. The same comparison at
the finitely many remarked split places follows by finite Howe
uniqueness. At all other places one can keep the original maps.
The support condition uses matrix rank in the polarization
specified by the chosen split embedding and is invariant under
invertible row and column changes. Therefore a pair matched with
the original maps is transported to an adapted matched pair after
rescaling one test by the product of these finitely many scalar
ratios. The two upper period functions are proportional also for
the original data, and a supported nonzero matched pair transports
back. This transport is needed only for the existence assertion
for an arbitrary marking. For the uniform recursion with fixed,
already adapted geometries use the direct calculation in the
first part of the proof. The upper period scalar may depend on
the chosen matching maps and on \(X\).
\end{proof}
\src{src/Development/WeilBound/ThetaLifts/LocalTheory/LocAOf.lean}{17}{The unprimed test comparison},
\src{src/Development/WeilBound/ThetaLifts/LocalTheory/LocBOf.lean}{17}{the primed test comparison}, and
\src{src/Development/WeilBound/ThetaLifts/LocalTheory/StepFacEq.lean}{85}{the quantifiers for a single descent constant}.
